\documentclass[tikz,border=4pt]{standalone}
\usepackage{ctex}
\usepackage{amsmath}
\usetikzlibrary{calc, arrows.meta}

\begin{document}
% thinking-toolkit/08 参数策略（含参讨论）
% 示意：关于x的方程 x^2+kx+1=0，按k的范围讨论根的情况
\begin{tikzpicture}[scale=1.0, thick]
  % 标题
  \node[font=\small, above] at (0,2.5) {含参方程 $x^2 + kx + 1 = 0$ 的讨论};

  % 判别式 Δ = k^2 - 4
  \node[font=\small] at (0,1.8) {判别式：$\Delta = k^2 - 4$};

  % 三种情形
  % Δ<0
  \draw[rounded corners=3pt, fill=pink!30, draw=black]
    (-5.5,1.2) rectangle (-2.5,0.0);
  \node[font=\footnotesize, align=center] at (-4,0.6)
    {$k^2 < 4$\\即 $-2<k<2$\\$\Delta<0$，无实根};

  % Δ=0
  \draw[rounded corners=3pt, fill=orange!30, draw=black]
    (-1.2,1.2) rectangle (1.2,0.0);
  \node[font=\footnotesize, align=center] at (0,0.6)
    {$k=\pm2$\\$\Delta=0$\\两相等实根};

  % Δ>0
  \draw[rounded corners=3pt, fill=green!20, draw=black]
    (2.5,1.2) rectangle (5.5,0.0);
  \node[font=\footnotesize, align=center] at (4,0.6)
    {$|k|>2$\\$\Delta>0$\\两不等实根};

  % 箭头从上方引下
  \draw[->,>=Stealth] (-0.8,1.7) -- (-4,1.25);
  \draw[->,>=Stealth] (0,1.7) -- (0,1.25);
  \draw[->,>=Stealth] (0.8,1.7) -- (4,1.25);

  % 数轴小图展示k的范围
  \draw[->] (-5.5,-0.8) -- (5.5,-0.8) node[right,font=\small] {$k$};
  \foreach \x/\lbl in {-2/$-2$, 0/$0$, 2/$2$}{
    \draw (\x,-0.7) -- (\x,-0.9) node[below,font=\footnotesize] {\lbl};
  }
  \fill[pink!30] (-2,-0.8) circle(3pt);
  \fill[pink!30] (2,-0.8) circle(3pt);
  \draw[very thick, pink!60] (-2,-0.8) -- (2,-0.8);
  \draw[very thick, green!60!black] (-5.5,-0.8) -- (-2,-0.8);
  \draw[very thick, green!60!black] (2,-0.8) -- (5.5,-0.8);
  \fill[orange!60] (-2,-0.8) circle(3pt);
  \fill[orange!60] (2,-0.8) circle(3pt);
  \node[above, font=\footnotesize, pink!60!black] at (0,-0.7) {无实根区};
  \node[above, font=\footnotesize, green!60!black] at (-4,-0.7) {两根区};
  \node[above, font=\footnotesize, green!60!black] at (4,-0.7) {两根区};
\end{tikzpicture}
\end{document}
